\section{Related Work}
\vspace{-5pt}
\noindent{\textbf{3D Gaussian Splatting for Novel View Synthesis.}} 
NeRF~\cite{mildenhall2021nerf} established a dominant paradigm for neural scene representation and sparked extensive subsequent research~\cite{martin2021nerf, barron2021mip, barron2022mip, yu2021plenoctrees, muller2022instant, fridovich2022plenoxels, park2021nerfies, barron2023zip, chen2022tensorf}, driving rapid progress in neural scene reconstruction. However, its per-ray volumetric rendering incurs high compute cost, motivating more efficient alternatives.
3DGS~\cite{kerbl20233dgs} mitigates this inefficiency by representing a scene with a set of 3D Gaussian primitives and rendering them through differentiable rasterization, enabling real-time rendering and faster optimization.
To achieve high fidelity with a compact representation, it further employs \textit{adaptive density control} (ADC), which periodically adds Gaussians in under-represented regions and prunes primitives with negligible contribution during iterative optimization.
This has inspired a line of work on more compact 3DGS representations, spanning both refinements of the ADC strategy~\cite{ye2024absgs, rota2024revising, zhang2024pixel, kim2024color, kheradmand20243d, cheng2024gaussianpro, zhang2024fregs, grubert2025improving, fang2024mini, mallick2024taming} 
and various compaction and pruning methods~\cite{lee2024compact, niedermayr2024compressed, papantonakis2024reducing, fan2024lightgaussian, girish2024eagles, chen2024hac, wang2024end, yang2024spectrally}.

Despite the advantages of 3DGS, it still has several practical limitations. It typically assumes dozens to hundreds of diverse input views for stable reconstruction, which can be impractical in real-world scenarios. This dense-view requirement has been partly addressed by recent studies on \textit{sparse-view} 3DGS~\cite{xiong2023sparsegs, zhang2024cor, zhu2024fsgs, li2024dngaussian, he2025see, kong2025generative}, which aim to reconstruct 3D scenes from only a few input images. 
Another major limitation is that 3DGS still requires iterative per-scene optimization, which remains a significant burden for practical deployment. To reduce this time-consuming optimization process, a variety of recent approaches~\cite{feng2025flashgs, zhao2024grendel, hollein20253dgslm, chen2025dashgaussian, wang2025grouptraining3dgs} have sought to accelerate the original 3DGS optimization pipeline through more efficient rasterization, parallelization, and improved optimization strategies. Among these approaches, \textit{feed-forward 3DGS} represents a particularly promising paradigm, as it amortizes iterative optimization into a single feed-forward pass, thereby enabling much faster 3D reconstruction.

% \cite{szymanowicz2024splatter, zheng2024gps, zhang2024gs, wewer2024latentsplat, szymanowicz2025flash3d, xu2025depthsplat, wang2024freesplat} 
\noindent{\textbf{Feed-Forward 3D Gaussian Splatting.}} Feed-forward 3DGS approaches~\cite{charatan2024pixelsplat, chen2024mvsplat, wang2024freesplat, chen2024pref3r, smart2024splatt3r, ye2024nopo, zhang2025flare, hong2024pf3plat, kang2025selfsplat, huang2025spfsplat, li2025vicasplat, jiang2025anysplat} have been proposed to alleviate the costly per-scene optimization of standard 3DGS. These methods are trained on large-scale datasets to learn strong priors, allowing them to predict 3D Gaussian representations in a single feed-forward pass without iterative optimization.
Consequently, they can reconstruct from sparse views while enabling real-time rendering and generalization to unseen scenes. Early generalizable feed-forward 3DGS methods~\cite{charatan2024pixelsplat, chen2024mvsplat, wang2024freesplat} typically assume calibrated multi-view inputs with known camera poses. Recent works~\cite{chen2024pref3r, smart2024splatt3r, ye2024nopo, zhang2025flare} relax this assumption by moving to pose-free settings. In addition, self-supervised pose-free approaches~\cite{hong2024pf3plat, kang2025selfsplat, huang2025spfsplat} further reduce reliance on pose annotations by learning from reconstruction consistency, with pose estimation integrated into the pipeline. More recently, uncalibrated formulations~\cite{li2025vicasplat, jiang2025anysplat} are enabling reconstruction without camera calibration.

Despite these advances in efficiency and robustness, existing feed-forward 3DGS methods largely rely on a \emph{uniform} output parameterization, in which a fixed number of Gaussians is allocated per pixel or spatial unit. As a result, the total Gaussian count becomes tightly coupled with the input resolution, rather than being \emph{adaptively} allocated according to scene complexity. This leads to redundant primitives in simple regions, while failing to sufficiently model geometrically complex regions, resulting in a suboptimal and non-compact representation under a limited Gaussian budget. In conventional optimization-based 3DGS, this issue has been mitigated through adaptive density control (ADC), which dynamically allocates Gaussians based on scene structure. However, such mechanisms rely on iterative per-scene optimization and are therefore not directly applicable to feed-forward pipelines. The recent feed-forward approach, AnySplat~\cite{jiang2025anysplat}, is able to control the Gaussian count via voxel granularity. However, the allocation remains spatially uniform, and adapting to different budgets typically requires retraining, limiting flexibility and compactness. In contrast, we introduce Gaussian-count controllable feed-forward 3DGS, which predicts a budget-aware densification score that enables non-uniform and spatially adaptive Gaussian allocation. This yields a more compact 3D representation under a controllable Gaussian budget.